\documentclass[a4paper,11pt]{article}
\usepackage[T2A]{fontenc}
\usepackage[utf8]{inputenc}
\usepackage[russian]{babel}
\usepackage{amsmath,amssymb,mathtools}
\usepackage{cmbright}
\renewcommand{\familydefault}{\sfdefault}
\usepackage{icomma}
\usepackage[a4paper,left=16mm,right=16mm,top=12mm,bottom=11mm]{geometry}
\usepackage{microtype}
\usepackage{multicol}
\usepackage{xcolor}
\usepackage{hyperref}
\definecolor{accent}{HTML}{15747A}
\definecolor{darkaccent}{HTML}{164C55}
\definecolor{textgray}{HTML}{283238}
\definecolor{muted}{HTML}{65747A}
\definecolor{rulegray}{HTML}{CADADC}
\color{textgray}
\hypersetup{hidelinks}
\setlength{\parindent}{0pt}
\setlength{\parskip}{0pt}
\setlength{\columnsep}{8mm}
\setlength{\abovedisplayskip}{1mm}
\setlength{\belowdisplayskip}{1mm}
\emergencystretch=2em
\newcommand{\sect}[1]{\par\vspace{2.2mm}{\normalsize\bfseries\color{darkaccent}#1}\par\vspace{0.3mm}{\color{rulegray}\rule{\linewidth}{0.45pt}}\par\vspace{0.8mm}}
\newcommand{\fml}[1]{\par\noindent{\normalsize\(\displaystyle #1\)}\par\vspace{2.8mm}}
\begin{document}
\thispagestyle{empty}

\hypersetup{pdftitle={Главные формулы. Екатерина Гнедкова},pdfauthor={Лёвин Артём Александрович}}
\begin{center}
{\small\bfseries\color{accent}УЧЕБНАЯ ПАМЯТКА \quad\textbullet\quad ОСНОВНЫЕ ФОРМУЛЫ}\par
\vspace{1mm}
{\fontsize{21}{25}\selectfont\bfseries\color{darkaccent} Проценты, вклады, кредиты и векторы\par}
\vspace{0.8mm}
{\small\color{muted}Екатерина Гнедкова\quad\textbullet\quad ЕГЭ · профильная математика\quad\textbullet\quad 12.10.2026}
\end{center}
\vspace{2mm}
\noindent{\color{accent}\rule{\linewidth}{0.9pt}}
\vspace{1.3mm}
{\bfseries\color{darkaccent}Главные принципы решения}\par
\vspace{1mm}
\begin{multicols}{2}
\raggedcolumns
\small
\noindent\textbullet\quad Процентную ставку всегда переводите в долю: r=p/100.\par\vspace{0.75mm}
\noindent\textbullet\quad Для каждого периода отделяйте начисление процентов от выплаты.\par\vspace{0.75mm}
\columnbreak
\noindent\textbullet\quad В смесях сохраняется масса чистого вещества.\par\vspace{0.75mm}
\noindent\textbullet\quad Финансовые схемы записывайте таблицей «долг — проценты — платёж».\par\vspace{0.75mm}
\end{multicols}
\vspace{1.1mm}
\noindent{\color{accent}\rule{\linewidth}{0.5pt}}
\vspace{0.6mm}
\begin{multicols}{2}
\raggedcolumns
\normalsize
\sect{Проценты и вклады}
\fml{k=1+\frac p{100}}
\fml{S_n=S_0k^n}
\fml{S_{i}=kS_{i-1}+A_i}
\sect{Кредиты}
\fml{B_i=kB_{i-1}-X_i}
\fml{B_0=S,\quad B_n=0}
\fml{Sk^n=X(1+k+\cdots+k^{n-1})}
\columnbreak
\sect{Дифференцированная схема}
\fml{P_i=\frac Sn+r\frac{S(n-i+1)}n}
\fml{r=\frac p{100},\quad i=1,\ldots,n}
\sect{Смеси и векторы}
\fml{c=\frac{m_{\text{в-ва}}}{m_{\text{р-ра}}}}
\fml{c_1m_1+c_2m_2=c(m_1+m_2)}
\fml{\overrightarrow{AB}=(x_B-x_A;\,y_B-y_A)}
\end{multicols}
\vfill
{\color{accent}\rule{\linewidth}{0.8pt}}\par\vspace{1.5mm}
{\small\color{muted}Лёвин Артём Александрович эксклюзивно для Екатерины Гнедковой}
\end{document}